\documentclass[13pt, a4paper]{extarticle} % Hỗ trợ cỡ chữ 13pt chuẩn BM18

% ----- CÀI ĐẶT FONT VÀ NGÔN NGỮ (DÙNG PDFLATEX) -----
\usepackage[utf8]{inputenc} % Hỗ trợ gõ tiếng Việt
\usepackage[T5]{fontenc}    % Bảng mã tiếng Việt
\usepackage{mathptmx}       % Font Times mổ phỏng cho pdfLaTeX
\usepackage[vietnamese]{babel}

% ----- CÀI ĐẶT KHỔ GIẤY VÀ LỀ -----
\usepackage{geometry}
\geometry{
    a4paper,
    left=3cm, right=2cm, top=2.5cm, bottom=2.5cm % Lề tiêu chuẩn BM18
}

% ----- CÀI ĐẶT ĐÁNH SỐ TRANG GIỮA BÊN TRÊN -----
\usepackage{fancyhdr}
\pagestyle{fancy}
\fancyhf{} % Xóa mọi định dạng header/footer cũ
\fancyhead[C]{\thepage} % Đánh số trang ở chính giữa phía trên
\renewcommand{\headrulewidth}{0pt} % Xóa đường kẻ ngang ở header

% ----- CÁC GÓI TOÁN HỌC VÀ HÌNH ẢNH DÀNH CHO VRPTW -----
\usepackage{amsmath, amssymb}
\usepackage{graphicx}
\usepackage{caption}
\usepackage{indentfirst} % Thụt lề đoạn đầu tiên

% ----- CÀI ĐẶT TRÍCH DẪN (AUTHOR-YEAR) -----
\usepackage[authoryear, round]{natbib}
\renewcommand{\bibname}{Tài liệu tham khảo}

\begin{document}

% ================= TRANG BÌA =================
\begin{titlepage}
    \thispagestyle{empty} % TẮT ĐÁNH SỐ TRANG CHO TRANG BÌA
    \begin{center}
        TỔNG LIÊN ĐOÀN LAO ĐỘNG VIỆT NAM \\
        TRƯỜNG ĐẠI HỌC TÔN ĐỨC THẮNG \\
        --------------------------------- \\
        \vspace{2.5cm}
            VERSION THỬ, COI NHƯ NHÁP, KHÔNG DÙNG CÁI NÀY ĐƯA CÔ COI, NĂN NỈ NĂN NỈ NĂN NỈ
            \vspace{1.5cm}
        
        Công trình Nghiên cứu khoa học sinh viên năm học 2025 - 2026 \\
        \vspace{1.5cm}
        
        ĐỀ TÀI \\
        \vspace{0.5cm}
        \textbf{\Large TỐI ƯU HÓA BÀI TOÁN VRPTW SỬ DỤNG MẠNG NƠ-RON HỌC TĂNG CƯỜNG SÂU KẾT HỢP TÌM KIẾM LÂN CẬN    LỚN THÍCH ỨNG (DDQN-ALNS)} 
        \vspace{1.5cm}
        \\ 
        \text{\Large Giải pháp phiên bản 9.5 - Hybrid Plateau Solver} \\
        \vspace{2.5cm}

        ĐƠN VỊ: KHOA CÔNG NGHỆ THÔNG TIN
        \vspace{2.5cm}
    \end{center}
      \begin{flushright}

    \end{flushright}
    \begin{flushleft}
        GIẢNG VIÊN HƯỚNG DẪN: \\
        TS. Hồ Thị Linh \\
        \vspace{0.8cm}
        SINH VIÊN THỰC HIỆN: \\
        1. Huỳnh Nhật Huy \\
        2. Nguyễn Nhật Huy \\
        3. ...... .... .... Trân

    \end{flushleft}
    
    \vfill
    \begin{center}
        TP. Hồ Chí Minh, tháng 4 năm 2026
    \end{center}
\end{titlepage}

\newpage
\setcounter{page}{1}

%========================================================== INTRODUCTION
\section*{GIỚI THIỆU VÀ ĐẶT VẤN ĐỀ}
\textbf{Bài toán Định tuyến xe trong cửa sổ thời gian (Vehicle Routing Problem - VRPTW)}. Mở rộng từ bài toán CVRP (Capacitated Vehicle Routing Problem), trong VRPTW, khách không chỉ yêu cầu một lượng hàng hoá nhất định, phương tiện không chỉ bị giới hạn bởi tải trọng vật lý mà còn chỉ định rõ một cửa sổ thời gian cụ thể (time window - TW), việc giao nhận hàng bắt buộc phải được tiến hành. Phương tiện đến sớm hơn thời điểm trong cửa sổ thời gian này sẽ phải dừng lại và chờ cho tới khi TW nhả; và phương tiện không được muộn hơn thời điểm kết thúc TW. \vspace{0.25cm}
 
 Tổng quát hóa của Bài toán Traveling Salesman Problem - TSP, tìm một tuyến duy nhất đi qua tất cả điểm trên bản đồ. Bổ sung TW cắt không gian tìm kiếm dữ dội, khiến việc tìm kiếm một giải pháp khả thi (feasible solution) đơn thuần đối với VRPTW - tức giải pháp không vi phạm ràng buộc nào về tải trọng hay thời gian. \vspace{0.25cm}

Thuật toán Quy hoạch tuyến tính - Dantzig \& Ramser (1959), nền móng VRP. Tiên phong cho các thuật toán quy hoạch tuyến tính sau này.Ứng dụng toán tối ưu vào logistics. Clarke \& Wright (1964) đề xuất thuật toán tiết kiệm (Savings algorithm), một phương pháp xấp xỉ tham lam (greedy heuristic), gán mỗi khách vào một tuyến trực tiếp từ trạm gốc (giải pháp hợp lệ), rồi lại ghép nối hai tuyến lại, trên giá trị "tiết kiệm" nhất ( $S_{ij} = c_{i0} + c_{0j} - c_{ij}$). Chứng minh giới hạn thực tế có thể được xử lý hiệu quả qua cơ chế heuristic đơn giản không cần sức tính toán lớn. \vspace{0.25cm}

Brian Kallehauge củng cố và phân tích sâu trong các luận án và tài liệu nghiên cứu chuyên sâu về chủ đề. Thuật ngữ VRPTW chính thức được chuẩn hóa, phổ biến rộng rãi và trở thành một nhánh nghiên cứu độc lập mạnh mẽ liền với các công trình kinh điển của Marius M. Solomon và Jacques Desrosiers vào cuối những năm 1980.\vspace{0.25cm}

Solomon \& Desrosiers (1987, 1988) đã công bố bộ dữ liệu chuẩn (benchmark) cho bài toán, tạo ra thước đo khách quan cho mọi thuật toán tối ưu từ đó đến nay, là "tiêu chuẩn vàng" ngành. Hệ thống hóa các chiến lược heuristic, thiết lập VRPTW như một trụ cột trung tâm trong lý thuyết vận trù học.\vspace{0.25cm}

Desrochers, Desrosiers \& Solomon (1992) dùng phương pháp Dantzig-Wolfe với Column Generation. Biến VRPTW thành hai bài toán phối hợp nhau: Bài toán Chủ (Master Problem): phân hoạch tập hợp (Set Partitioning Problem). Chọn một tập các tuyến đường hợp lệ từ một pool các tuyến đường khả thi, sao cho tổng phí vận hành nhỏ nhất và mọi khách được phủ đúng một lần và Bài toán Con (Pricing Subproblem): nhận các giá trị đối ngẫu (dual values) phản hồi từ Master và cố gắng tìm kiếm các tuyến đường mới có chi phí rút gọn âm (negative reduced cost). Nếu thấy, tuyến này sẽ được thêm (sinh cột - column generation) vào Master để cải thiện nghiệm. Đối với VRPTW, bài toán con này là Bài toán Đường đi Ngắn nhất Sơ cấp có Ràng buộc Tài nguyên (Elementary Shortest Path Problem with Resource Constraints - ESPPRC). Các tài nguyên như tải trọng tích lũy, thời gian trôi và lượt ghé được liên tục cập nhật và kiểm soát dọc theo các tuyến tiềm năng.\\
Tuy nhiên, Column Generation độc lập chỉ giải quyết được dạng nới lỏng tuyến tính (linear relaxation) của Master. Các kết quả có thể có phân số (kiểu 0.5 chiếc xe để đi một tuyến đường) vô lý. Do đó, cấu trúc sinh cột bắt buộc phải được lồng ghép vào bên trong cây duyệt Branch-and-Bound (Nhánh \& Cận), sinh Branch-and-Price.\vspace{0.25cm}

Kohl \& Madsen (1997) mở rộng bằng thuật toán kết hợp ứng dụng nới lỏng Lagrangian (Lagrangian relaxation).\vspace{0.5cm}

Sau đó, Kohl, Desrosiers, Madsen, Solomon \& Soumis (1999) kết hợp thêm các bất phương trình cắt (cutting planes) để thu hẹp không gian đa diện nghiệm, đánh dấu sự ra đời của khung thuật toán Branch-and-Cut-and-Price.Sự kết hợp ba mũi nhọn này - vừa liên tục sinh các cột tuyến đường mới (Price) bằng quy hoạch động, vừa thêm các mặt cắt toán học để thu hẹp không gian đa diện (Cut), vừa rẽ nhánh để ép các biến phân số trở về số nguyên (Branch) - đã tạo nên chuẩn mực "State-of-the-Art" vô đối trong nhánh Exact Methods. \\
Nhờ các kỹ thuật này, giới hàn lâm dần có thể giải quyết dứt điểm và cung cấp chứng minh độ tối ưu cho các bộ dữ liệu Solomon, thiết lập các kỷ lục. Dù cực kỳ mạnh về mặt chứng minh toán học, hạn chế của Exact Methods là thời gian tính. Với các bài toán có quy mô vượt quá 200 khách hoặc những bài toán có cấu trúc ranh giới thời gian lỏng (nhóm R2, C2, RC2 Solomon), các công cụ phân rã toán học phải mất hàng ngày, tuần, hoặc không thể hội tụ, quá chậm để áp dụng trong hoạt động điều phối thực tiễn. Điều này đã mở đường cho một hướng đi song song định hình ngành logistics: Kỷ nguyên Heuristics và Metaheuristics.\vspace{0.5cm}

Sự phức tạp cao và yêu cầu khắt khe về thời gian giải quyết nhanh gọn cho các kịch bản thực tiễn (logistics hàng ngày) đã nâng tầm quan trọng của các phương pháp xấp xỉ (heuristics và metaheuristics). Những phương pháp này chấp nhận hy sinh việc đảm bảo tìm được kết quả tối ưu tuyệt đối, nhưng bù lại chúng cung cấp các cấu trúc tuyến đường có chất lượng cực kỳ cao (near-optimal) trong những khoảng thời gian khả thi (vài giây đến vài phút). Theo thời gian, những phương pháp này đã tiến hóa từ các quy tắc đơn giản thành những hệ thống thông minh tinh vi.\vspace{0.5cm}

Trong giai đoạn đầu của thập niên 1990 và đầu 2000, vô số thuật toán metaheuristic đã được thử nghiệm và tinh chỉnh trên VRPTW. Các phương pháp ban đầu tập trung vào Tìm kiếm Lân cận Cục bộ (Local Search), trong đó một giải pháp khởi tạo (initial solution) được liên tục tinh chỉnh thông qua các toán tử hoán đổi đỉnh hoặc cạnh đơn giản như Or-opt.\\
Or-opt: Or (1976), cho phép thuật toán di dời một chuỗi các đỉnh liên tiếp sang một vị trí khác trên tuyến đường mà không làm biến đổi định hướng của tuyến gốc.Theo sau đó là sự áp dụng của các chiến lược định hướng toàn cục mạnh mẽ hơn:\\
Tabu Search: Gendreau et al. (1999) đã đề xuất sử dụng Tabu Search để ngăn chặn thuật toán bị kẹt tại các cực tiểu cục bộ. Phương pháp này sử dụng một danh sách bộ nhớ (Tabu list) để cấm thuật toán quay lại các giải pháp vừa được thăm trong một số chu kỳ nhất định, buộc quá trình tìm kiếm phải khám phá các vùng không gian nghiệm mới.\\
Ant Colony Optimization: Mô phỏng hành vi tìm thức ăn của loài kiến, nơi các "con kiến" xây dựng tuyến đường dựa trên xác suất được tinh chỉnh bởi nồng độ pheromone (được tích lũy từ các tuyến đường hiệu quả) và thông tin kinh nghiệm (heuristic information). Gambardella et al. (1999) đã phát triển mô hình MACS-VRPTW, chứng minh rằng mở rộng MACS là một trong những cách tiếp cận hiệu quả và tốc độ cao nhất thời bấy giờ cho cả biến thể định tuyến động và tĩnh có cửa sổ thời gian.\\

Evolutionary Algorithms (EA) \& Genetic Algorithms (GA): Thuật toán di truyền mô phỏng sự tiến hóa của sinh học thông qua quá trình lai ghép (crossover) và đột biến (mutation). Tuy nhiên, các thuật toán GA truyền thống gặp khó khăn nghiêm trọng khi xử lý ràng buộc thời gian khắt khe, vì việc lai ghép hai tuyến đường tốt thường tạo ra các thế hệ con không khả thi. Điều này dẫn đến sự ra đời của các "Thuật toán Memetic" (Memetic Algorithms) - sự lai tạo giữa GA và các bộ tìm kiếm cục bộ mạnh mẽ.\\
Để thúc đẩy việc đánh giá giới hạn của các thuật toán metaheuristic này trên những mạng lưới siêu lớn, Gehring và Homberger (1999) đã giới thiệu một bộ dữ liệu mở rộng (Extended Benchmark) bao gồm các bản đồ mô phỏng lên tới 200, 400, 600, 800 và 1000 khách hàng. Khối lượng tính toán khổng lồ từ các tập dữ liệu này đặt ra thách thức cực lớn mà chỉ những kiến trúc thuật toán xuất sắc nhất mới có thể xử lý, loại bỏ hoàn toàn khả năng can thiệp của các exact methods.\vspace{0.5cm}

Một trong những cột mốc vĩ đại nhất, định hình lại hoàn toàn cục diện state-of-the-art của các thuật toán metaheuristics là sự xuất hiện của khung thuật toán Adaptive Large Neighborhood Search (ALNS). Được Pisinger và Ropke giới thiệu vào năm 2006 và mở rộng cho VRPTW vào năm 2007, ALNS đã phá vỡ mọi giới hạn trước đó. Khác với các thuật toán tìm kiếm cục bộ chỉ tinh chỉnh những cạnh nhỏ lẻ và dễ bị kẹt, ALNS can thiệp ở quy mô không gian cực lớn thông qua cơ chế "Phá hủy và Xây dựng lại" (Ruin and Recreate).Đặc tính Kỹ thuậtCơ chế hoạt động trong ALNSTác động tối ưu hóaToán tử Phá hủy (Destroy Operators)Tước bỏ một lượng lớn khách hàng (thường từ 10% đến 40%) khỏi cấu trúc tuyến đường hiện tại. Các chiến lược phổ biến bao gồm Shaw Removal (xóa các khách hàng có sự tương đồng mạnh về vị trí và cửa sổ thời gian), Worst Removal (xóa các nút gây ra gia tăng chi phí cao nhất), hoặc Random Removal để tạo sự ngẫu nhiên.Phá vỡ triệt để các cấu trúc cực tiểu cục bộ cứng đầu, giải phóng các nút để chúng có cơ hội được gán vào các tuyến đường tối ưu hơn.Toán tử Sửa chữa (Repair Operators)Chèn các khách hàng vừa bị xóa trở lại mạng lưới bằng các chiến thuật thông minh như Greedy Insertion (chèn vào nơi tốn ít chi phí nhất) hoặc Regret-based Insertion (chèn khách hàng vào vị trí dựa trên độ hối tiếc - chênh lệch chi phí giữa vị trí tốt nhất và tốt nhì).Tái thiết lập các tuyến đường mới tinh vi hơn, thường dẫn đến việc loại bỏ hoàn toàn một số phương tiện (giảm số xe).Cơ chế Thích nghi (Adaptive Mechanism)Trọng số xác suất sử dụng mỗi toán tử phá hủy và sửa chữa được điều chỉnh động dựa trên hiệu suất thành công của chúng trong các vòng lặp trước đó (sử dụng vòng quay roulette).Thuật toán tự học hỏi và điều chỉnh chiến lược tùy thuộc vào đặc thù cấu trúc của từng instance dữ liệu cụ thể.Ngoài ra, ALNS áp dụng tiêu chí chấp nhận nghiệm của thuật toán Mô phỏng Luyện kim (Simulated Annealing - SA) để cho phép chấp nhận các nghiệm xấu hơn ở giai đoạn đầu với xác suất giảm dần, từ đó bảo đảm khả năng thoát khỏi các tối ưu cục bộ. ALNS đã trở thành hệ xấp xỉ tiêu chuẩn thống trị giải quyết VRPTW trong nhiều năm liên tiếp, duy trì vị thế state-of-the-art và liên tục được lấy làm cơ sở so sánh cho mọi thuật toán mới.

Mặc dù ALNS vô cùng xuất sắc, cấu trúc giải thuật tiến hóa bắt đầu giành lại ngôi vương từ năm 2012 với sự xuất hiện của kiến trúc Hybrid Genetic Search (HGS). Được phát triển bởi Thibaut Vidal, Teodor Gabriel Crainic, Michel Gendreau và Christian Prins, HGS đã thiết lập một kỷ lục hoàn toàn mới về độ chính xác và tính nhất quán, không chỉ cho VRPTW mà cho hầu hết các biến thể VRP trên toàn cầu. Đáng kinh ngạc là đến tận những năm 2024-2026, các biến thể mở rộng từ nền tảng HGS vẫn đứng vững ở vị thế thuật toán metaheuristic cổ điển tiên tiến nhất (state-of-the-art), vượt qua sự bào mòn của thời gian.6.1. Ba Trụ cột Thiết kế Đột phá của Hệ sinh thái HGSSự xuất sắc vô tiền khoáng hậu của HGS được cấu thành từ ba cơ chế cốt lõi được kết hợp một cách hoàn hảo :Quản lý Đa dạng Thích ứng (Adaptive Diversity Management): Trong các thuật toán di truyền (GA) truyền thống, quần thể thường gặp phải vấn đề hội tụ quá sớm (premature convergence), dẫn đến tình trạng các cá thể (nhiễm sắc thể tuyến đường) trở nên quá giống nhau, làm mất đi khả năng khám phá. HGS giải quyết triệt để vấn đề này bằng cách đánh giá chất lượng (fitness) của mỗi cá thể không chỉ thông qua hàm mục tiêu (chi phí và khoảng cách) mà còn thông qua độ đóng góp đa dạng (contribution to diversity) của cá thể đó so với quần thể hiện hành. Thuật toán thưởng cho các cá thể có cấu trúc di truyền độc đáo, duy trì nguyên lý "sống sót của sự khác biệt", qua đó bảo toàn các gen quy hoạch tuyến đường tiềm năng trong suốt hàng nghìn thế hệ lai ghép.Thuật toán rã chuỗi trong thời gian tuyến tính (Split Algorithm $O(n)$): Vidal và cộng sự tận dụng cơ chế biểu diễn nghiệm cực kỳ thông minh dưới dạng một chuỗi lộ trình khổng lồ ("Giant-Tour") đi qua tất cả khách hàng, hoàn toàn bỏ qua mọi ràng buộc về phương tiện và sức chứa. Để lượng giá chuỗi này, HGS áp dụng một thuật toán chia rẽ (Split algorithm) hiệu suất cực cao hoạt động trong thời gian $O(n)$, tối ưu hóa việc cắt chuỗi khổng lồ này thành các tuyến đường hợp lệ phân bổ cho từng xe dựa trên các ràng buộc. Việc tối ưu toán tử Split này giúp HGS thực hiện hàng triệu lần lai ghép trong vài giây.Thăm dò Không gian Bất khả thi có Phạt (Penalized Infeasible Search): Đây là chiến lược mang tính chất bước ngoặt quyết định sự thành bại đối với VRPTW. Thay vì bác bỏ hoàn toàn các giải pháp vi phạm ràng buộc như các thuật toán trước đó, HGS cố tình cho phép thuật toán tạm thời vượt rào các giới hạn về sức chứa và cửa sổ thời gian trong quá trình tìm kiếm cục bộ. Thông qua kỹ thuật "Time Warp" (bẻ cong thời gian), các xe được phép "đi ngược thời gian" để đến kịp cửa sổ cho phép, nhưng hệ thống sẽ cộng dồn một mức hình phạt (penalty factor) vào hàm mục tiêu. Mức hình phạt này được điều chỉnh động: nếu quần thể có quá nhiều nghiệm vi phạm, hình phạt tăng lên để đẩy chúng về vùng khả thi; nếu quần thể quá an toàn, hình phạt giảm xuống để khuyến khích khám phá. Việc đi xuyên qua ranh giới bất khả thi giúp HGS tìm ra các "đường hầm" (tunnels) kết nối các vùng tối ưu cục bộ bị cô lập trong không gian tìm kiếm, dẫn tới việc tìm thấy các giải pháp vượt trội chưa từng có.6.2. Mã nguồn mở và Uy lực của toán tử SWAP*Vào năm 2020 (và xuất bản chính thức trên tập san uy tín vào năm 2022), Thibaut Vidal đã tiến hành tối ưu hóa triệt để cấu trúc mã và công bố một phiên bản HGS mã nguồn mở dạng C++ chuyên biệt cho CVRP và VRPTW. Động thái này đã tạo ra một cơn địa chấn trong cộng đồng nghiên cứu, cung cấp một hệ quy chiếu mã nguồn đỉnh cao cho mọi so sánh.Phiên bản này tích hợp thêm một toán tử thăm dò vùng lân cận cực kỳ uy lực mang tên SWAP* (phát triển từ thuật toán của các nhà nghiên cứu đi trước). Khác với SWAP thông thường chỉ tráo đổi hai điểm ở cùng vị trí vật lý, toán tử SWAP* cho phép trao đổi hai khách hàng giữa các tuyến đường khác nhau mà không yêu cầu việc chèn khách hàng vào đúng vị trí cũ; thay vào đó, nó tính toán và rà soát vị trí chèn lý tưởng nhất dọc theo toàn bộ tuyến đường mục tiêu để tối ưu hóa sự tương thích không gian - thời gian. Bằng sự kết hợp giữa kiến trúc di truyền, quản lý đa dạng và không gian lân cận SWAP*, HGS đã chính thức thiết lập hàng loạt các Kỷ lục Giải pháp Tốt nhất Mới (New BKS) trên hầu hết mọi tập dữ liệu Solomon 100 và Homberger 1000 khách hàng, thống trị bảng xếp hạng trong nhiều năm.7. Kỷ nguyên Trí tuệ Nhân tạo và Định tuyến Bằng Học sâu (2024 - 2026)Mặc dù HGS đạt được những giới hạn tối ưu phi thường, các phương pháp metaheuristic truyền thống vẫn gặp phải rào cản nghiêm trọng: chúng yêu cầu thiết kế thủ công các quy tắc heuristic đòi hỏi chuyên môn cao, tính toán chuyên sâu và thời gian chạy mất từ hàng chục phút đến hàng giờ đối với các bộ dữ liệu lớn để hội tụ. Trong bối cảnh logistics thời gian thực, điều này là không thể chấp nhận.
Từ năm 2024 đến 2026, lĩnh vực Tối ưu hóa Tổ hợp (Combinatorial Optimization) chứng kiến một cuộc cách mạng sâu sắc khi giới học thuật và công nghiệp tích hợp Trí tuệ nhân tạo (AI), Học máy (ML), và đặc biệt là Học tăng cường sâu (Deep Reinforcement Learning - DRL) để giải quyết các hạn chế về thời gian thực. DRL cho phép các tác tử (agents) tự học ra các chiến lược tối ưu để giải quyết trực tiếp bài toán thông qua suy diễn đồ thị chỉ trong vài mili-giây (inference time), mở ra khả năng tái định tuyến liên tục (dynamic routing).7.1. Khung Học tăng cường sâu (DRL) tích hợp Mạng Neural Đồ thịĐể mô hình hóa phương pháp tiếp cận DRL cho VRPTW, quá trình định tuyến được quy về một dạng Quá trình Ra quyết định Markov (Markov Decision Process - MDP). Tại mỗi bước ra quyết định $t$, trạng thái mạng lưới (bao gồm tọa độ khách hàng, thời gian sẵn sàng, nhu cầu bưu kiện, dung lượng phương tiện hiện tại) được biểu diễn thành các tensor trạng thái. Trí tuệ nhân tạo sẽ huấn luyện một tác tử để lựa chọn nút tiếp theo cần đến nhằm tối đa hóa phần thưởng kỳ vọng (tương đương với giảm thiểu khoảng cách di chuyển và tránh hình phạt vi phạm thời gian).Những phương pháp state-of-the-art (SOTA) trong năm 2025-2026 thiết lập tiêu chuẩn mới bằng cách kết hợp các cơ chế Mạng Neural Đồ thị (GNN - Graph Neural Networks) cực kỳ tinh vi:Kiến trúc Encoder-Decoder (Bộ mã hóa - Bộ giải mã) lai: Các nghiên cứu DRL SOTA áp dụng kiến trúc mạng chú ý đồ thị có cổng tiên tiến (Gated Graph Attention Network - GATv2) cho bộ mã hóa (Encoder) để trích xuất một cách động các mối tương quan không gian phức tạp và độ khắt khe về thời gian giữa các nút khách hàng liền kề.Giải mã nhận thức ngữ cảnh (Dynamic Context-aware Decoding): Bộ giải mã sử dụng Cơ chế Chú ý Đa đầu (Multi-head Attention - MHA) kết hợp với mạng chuỗi thời gian LSTM (Long Short-Term Memory) để bảo tồn sự phụ thuộc về mặt thời gian (temporal dependencies) trong quá trình chọn tuần tự từng nút, giúp đánh giá và dự báo chính xác khả năng tuân thủ ranh giới thời gian cứng ở các cụm nút tiếp theo trong tương lai.Cơ chế tối ưu hóa chính sách và Action Masking: Phương pháp Tối ưu hóa Chính sách Cận kề (Proximal Policy Optimization - PPO) được sử dụng để duy trì tính ổn định của gradient trong quá trình đào tạo. Trong pha suy diễn, các hành động (chọn khách hàng) vi phạm sức chứa hoặc dẫn đến đến muộn không thể cứu vãn sẽ bị thuật toán loại bỏ hoàn toàn khỏi phân phối xác suất thông qua cơ chế Action Masking. Các hàm phần thưởng tổng hợp (composite reward functions) bao gồm hình phạt vi phạm ràng buộc (constraint-aware) cũng được tinh chỉnh để định hướng tác tử.Các nghiên cứu thực nghiệm quy mô lớn trên bộ dữ liệu ngẫu nhiên, tập dữ liệu Solomon chuẩn mực, và tập dữ liệu logistics phân phối thực tế của doanh nghiệp cho thấy: khung DRL lai này có khả năng làm giảm chi phí vận chuyển từ 2-10% so với các phương pháp DRL sơ khai, đạt tỷ lệ thỏa mãn ràng buộc và khả thi nghiệm vượt qua mốc 99%, trong khi thời gian thực thi (inference) nhanh hơn hàng trăm lần so với các hệ Branch-and-Price hoặc HGS khi chạy suy diễn trực tuyến.7.2. Các thuật toán định tuyến động: Sự cạnh tranh giữa DDQN và Mạng TransformerBên cạnh PPO, các mạng Q-learning chuyên sâu như Double Deep Q-Network (DDQN) cũng được triển khai sâu rộng cho những môi trường phân phối hàng hóa biến động mạnh (dynamic VRPTW), nơi đơn đặt hàng, thời gian nhận, hoặc tình trạng giao thông có tính chất bất định và liên tục thay đổi. Khác với DRL dựa trên chính sách (policy-based) như PPO thường đòi hỏi cập nhật không gian hành động liên tục, phương pháp tiếp cận dựa trên giá trị (value-based) như DDQN đánh giá trực tiếp các hành động khả thi rời rạc, cung cấp cơ chế ra quyết định nhẹ nhàng hơn về mặt tính toán cho định tuyến trực tuyến. Đặc biệt, cấu trúc Double (kép) của DDQN giúp loại bỏ hoàn toàn sự ước lượng sai lệch tăng cường (overestimation bias) – một nhược điểm chí mạng thường có ở DRL cổ điển (DQN) khi các giá trị Q khởi tạo quá cao làm lệch hướng huấn luyện.Một số cấu trúc DDQN tối tân (ví dụ: mô hình dự đoán cho Dada Express) đã tích hợp các hàm phần thưởng kép (dual-reward) sử dụng cả Tỷ lệ Khớp hoàn toàn (Complete Match Rate - CMR) cho sự căn chỉnh tuyến đường nghiêm ngặt và Khoảng cách Sửa đổi (Edit Distance - ED) làm phản hồi gia tăng. Cấu trúc này còn áp dụng chiến lược Experience Replay ưu tiên (APER) kết hợp kỹ thuật phân tích lộ trình phi tuyến nhằm tái tạo chính xác hành vi sở thích tuyến đường thực tế của tài xế (rider-specific routing preferences) - điều mà các mô hình Toán học thuần túy không bao giờ làm được.Tuy nhiên, trong quá trình đo kiểm benchmarking nghiêm ngặt của năm 2025 về khả năng nắm bắt liên kết mạng lưới (Link Prediction và Node dependencies), các kiến trúc sử dụng Mạng Transformer đang thể hiện sức mạnh áp đảo đối với việc tiếp thu ngữ cảnh đồ thị tuyến đường khổng lồ so với các mạng GNN hay GAT cơ bản. Khả năng tự chú ý (Self-Attention) toàn cục của Transformer cho phép nhận diện tầm nhìn chiến lược qua toàn bộ bản đồ, hoàn toàn bỏ qua sự phụ thuộc cục bộ của các nút liền kề trong kiến trúc GNN truyền thống. Sự cạnh tranh giữa DDQN, GATv2 và Transformer đang định hình bức tranh nghiên cứu thuật toán giao nhận thời gian thực.Phương pháp DRLƯu điểm cốt lõi trong VRPTWNhược điểm/Thách thứcPPO với GATv2/LSTMĐộ ổn định huấn luyện cực cao, mô hình hóa tốt độ tương quan không gian-thời gian cục bộ.Đòi hỏi tham số lớn, khả năng bắt tín hiệu toàn cục (global dependencies) trên bản đồ lớn kém hơn Transformer.Double DQN (DDQN)Triệt tiêu sai số đánh giá vượt mức (overestimation bias), hoạt động cực tốt trong điều kiện động (dynamic routing) thời gian thực.Khó khăn khi không gian hành động (số lượng khách hàng) trở nên quá lớn, tốn tài nguyên bộ nhớ đệm (Replay Buffer).Kiến trúc TransformerTầm nhìn đồ thị toàn cục thông qua Self-Attention, dễ dàng nắm bắt sự tương quan giữa 2 nút ở cách xa nhau trên bản đồ.Yêu cầu tài nguyên huấn luyện (GPU VRAM) khổng lồ, chi phí tính toán tăng bậc hai ($O(n^2)$) theo số nút.

Một bước ngoặt không tưởng của tiến trình tối ưu hóa VRP trong các năm 2025-2026 là sự can thiệp của Mô hình Ngôn ngữ Lớn (Large Language Models - LLMs) vào lĩnh vực cốt lõi nhất của Operations Research: tự động thiết kế các toán tử thuật toán (Automated Heuristics Discovery).Trong nhiều thập kỷ, các thuật toán (như Clarke-Wright, Or-opt, SWAP*) đều là những di sản trí tuệ do sự sáng tạo thủ công của các nhà toán học. Tuy nhiên, theo định hướng của các hệ quy chiếu học thuật tối tân như cuộc thi GECCO 2026 chuyên đề "Machine Learning Assisted Evolutionary Computation" , giới hàn lâm đã giới thiệu các mô hình AI Agent có khả năng thay thế con người trong tác vụ này.Đại diện tiêu biểu nhất là hệ thống VRPAgent. VRPAgent tích hợp một bộ tạo mã từ LLM (sử dụng các mô hình như Gemini 2.5 Flash, Qwen3, hoặc Llama 3.3) vào bên trong hệ sinh thái metaheuristic LNS (Large Neighborhood Search). Thông qua một quy trình tiến hóa di truyền trên "không gian mã nguồn" (Code Space), LLM sinh ra những đoạn mã Python chứa các hàm toán tử heuristic phá hủy/sửa chữa hoàn toàn mới, sáng tạo và mang tính đặc thù cho các dữ liệu VRPTW đầu vào. Mã nguồn này ngay lập tức được biên dịch, chạy thử trên CPU để lấy điểm đánh giá fitness, và sử dụng cơ chế bảo tồn tinh hoa (elitism) để phản hồi lại LLM làm bệ phóng cho việc viết lại đoạn mã tốt hơn trong thế hệ tiếp theo. Các kết quả đánh giá thực chứng trên bộ dữ liệu Solomon và CVRP cho thấy VRPAgent tự động khám phá ra các chiến lược vượt trội hơn cả mã nguồn thiết kế thủ công từ chuyên gia, đồng thời thiết lập SOTA mới trong lớp thuật toán học máy giải quyết VRP mà chỉ yêu cầu sử dụng một lõi CPU đơn lẻ ở khâu thực thi.Hệ thống mang tính đột phá về cấu trúc tư duy hơn nữa là Metacognitive Evolutionary Programming (MEP) (thường được nhắc tới thông qua nền tảng PyVRP+ do nhóm nghiên cứu Malik, Zhou et al. phát triển). Phương pháp này ép buộc LLM không chỉ sinh mã như một cỗ máy đột biến hộp đen thông thường, mà phải tiến hành một chu trình "Lý luận - Hành động - Suy ngẫm" (Reason-Act-Reflect) có nhận thức siêu việt. Trước khi tạo mã mới, LLM sẽ phân tích log thất bại tại sao toán tử hiện tại bị kẹt ở điểm tối ưu cục bộ, từ đó đặt ra các giả thuyết thiết kế khoa học, giải quyết bài toán cân bằng giữa thăm dò và khai thác (exploration-exploitation trade-off), rồi mới viết ra thuật toán lân cận mới dựa trên kiến thức miền sâu sắc. Việc ứng dụng MEP để tiến hóa các toán tử lõi trong kiến trúc Hybrid Genetic Search (HGS) cổ điển đã giúp giảm tới 45% thời gian chạy tính toán trong khi tối ưu chất lượng nghiệm thêm 2.70% trên các kịch bản khó nhất, minh chứng sức mạnh vô tận của LLM-driven optimization.

9.1. Định hướng lượng tử hóa thuật toán VRP (Quantum Computing)Trong nỗ lực nhằm tìm ra cách xuyên phá bức tường bùng nổ tổ hợp, giai đoạn từ 2024 đến 2026 đã ghi nhận những chứng minh thực tế đầu tiên mang tính bước ngoặt về việc ứng dụng điện toán lượng tử vào bài toán VRPTW.Sự trỗi dậy của máy tính ủ lượng tử (Quantum Annealing) được dẫn dắt bởi hệ thống Advantage2 của công ty D-Wave. Bằng cách ánh xạ mô hình toán học của VRPTW (với các biến quyết định nhị phân) sang các phương trình năng lượng biểu diễn Hamilton (dưới định dạng QUBO - Quadratic Unconstrained Binary Optimization) kết hợp mô hình học máy lai lượng tử-cổ điển, máy tính D-Wave có thể kiểm tra một không gian năng lượng tổng thể cực đại chỉ trong nháy mắt thay vì duyệt qua hàng tỷ biến thể bằng CPU truyền thống. Các thử nghiệm thực tế tại châu Âu vào đầu năm 2025 đã ứng dụng Quantum Annealing để tối ưu hóa điều phối giao thông công cộng, chứng minh việc giảm chi phí đáng kinh ngạc trong thời gian thực.Bên cạnh máy ủ lượng tử, các nền tảng lượng tử dựa trên cổng logic (Gate-based Quantum Computers) cũng đưa ra các khung Thuật toán Tối ưu hóa Xấp xỉ Lượng tử (QAOA - Quantum Approximate Optimization Algorithm). Nhược điểm truyền thống của QAOA trên các bài toán có nhiều giới hạn chặt chẽ về mặt không gian và thời gian như VRPTW là thuật toán lượng tử thường lạc lối và lãng phí tài nguyên vào các không gian nghiệm không hợp lệ (infeasible space). Một nghiên cứu tiên phong (Lei et al. 2025) đã khắc phục điểm yếu này bằng một kỹ thuật "Constraint-aware QAOA" tích hợp hệ thống bộ trộn (mixer) XY-X lai. Cấu trúc mixer này bảo tồn một phần các tuyến đường đã hợp lệ trong khi tiến hành khám phá các tuyến đường mới qua cổng lượng tử, giúp nâng tỷ lệ tìm thấy giải pháp khả thi lên đến 99.7%, mở khóa sự khả thi của định tuyến chuỗi cung ứng trên lượng tử.Song song đó, một giao điểm nghiên cứu cực kỳ hứa hẹn là sự hình thành của kiến trúc mạng Quantum Graph Attention Network (Q-GAT). Mô hình này thay thế trực tiếp các mạng Perceptron Đa lớp (MLPs) khổng lồ và tiêu tốn bộ nhớ trong các pha trích xuất thông tin (readout) của mạng học sâu DRL bằng các Mạch Lượng tử Tham số (Parameterized Quantum Circuits - PQCs). Quá trình này giúp mô hình giảm hơn 50% số lượng tham số học máy trong khi vẫn giữ nguyên năng lực phân tích biểu diễn cấu trúc đồ thị chú ý. Khi kết hợp mã hóa lượng tử với chính sách PPO trên các bộ giải mã tham lam và ngẫu nhiên, Q-GAT đạt trạng thái hội tụ nhanh hơn và giảm 5% chi phí lộ trình so với các đường cơ sở GAT thông thường (như GATv2 cổ điển), chứng minh tiềm năng là bộ giải mã lý tưởng cho logistics quy mô lớn trên thiết bị điện toán biên.9.2. Sự trỗi dậy của tối ưu hóa tuyến tính trên GPU (NVIDIA cuOpt)Nếu điện toán lượng tử mang tính định hướng trung và dài hạn (kỳ vọng tích hợp hoàn toàn vào quy trình làm việc sau 2026-2032) , thì việc sử dụng tài nguyên xử lý song song khổng lồ của GPU (Graphics Processing Unit) đã cung cấp khả năng tăng tốc độ tuyến tính giải quyết bài toán ngay lập tức. Công cụ tối ưu hóa quyết định nguồn mở NVIDIA cuOpt đại diện cho nỗ lực thu hẹp khoảng cách giữa các thư viện giải thuật truyền thống và năng lực xử lý phân luồng cấp cao của kiến trúc vi xử lý đồ họa.Thay vì dựa vào thuật toán truyền thống Simplex hoạt động trên kiến trúc CPU tuần tự (vốn rất chậm khi giải quyết Linear Programming cho VRP), cuOpt tận dụng cơ chế Primal-Dual Linear Programming (PDLP) hiện đại. PDLP là một phương pháp bậc nhất (First-Order Method) tính toán xấp xỉ đạo hàm của bài toán và giảm thiểu vi phạm ranh giới thời gian bằng các thao tác tính toán ma trận song song khối lượng lớn. Đối với các hệ bài toán hỗn hợp MIP (Mixed Integer Programming) và VRPTW có hàng triệu biến số cấu thành, sự hỗ trợ từ hệ sinh thái CUDA và băng thông đọc ghi cực kỳ khổng lồ từ các thế hệ GPU mới nhất của NVIDIA (như HGX B100 đạt tới 8 TB/s) đã giúp trình giải cuOpt đạt tốc độ tính toán nhanh hơn gấp 5.000 lần so với các trình giải trên CPU phổ thông.Quy trình giải quyết (solver logic) của cuOpt sử dụng kiến trúc kết hợp GPU/CPU tinh vi: các phương pháp tìm kiếm nghiệm xấp xỉ nguyên thủy (primal heuristics) như local search, feasibility pump, và feasibility jump hoạt động trực tiếp với tính chất song song hàng loạt trên GPU để cải thiện các ràng buộc cận (primal bound), trong khi các thao tác xác định giới hạn đối ngẫu phức tạp Branch-and-bound (duyệt cây nhị phân) được chuyển giao cho CPU xử lý. Các giải pháp nguyên hợp lệ được chia sẻ liên tục giữa cả hai thuật toán giúp tăng tốc độ chốt nghiệm.10. Các Tiêu chuẩn Đánh giá Benchmark và Hệ thống Cuộc thi Quốc tếSự tiến hóa về toán học thuật toán và ứng dụng mạng thần kinh chỉ có ý nghĩa học thuật nếu có các bộ dữ liệu tiêu chuẩn (Benchmarks) minh bạch làm nền tảng kiểm định năng lực. Mọi kiến trúc thuật toán từ những năm 1980 cho đến hiện tại đều xoay quanh việc đánh bại các giới hạn trên tập hợp Benchmark tiêu chuẩn này.10.1. Bộ dữ liệu Solomon và Tập mở rộng Gehring & HombergerNhư đã phân tích, tập dữ liệu Solomon gồm các bài toán 100 khách hàng (công bố năm 1987) đã đóng vai trò lịch sử vĩ đại. Đến tận năm 2026, nó vẫn đóng vai trò là một tài liệu chuẩn mực bắt buộc phải vượt qua đối với bất kỳ bài báo nào khảo sát VRPTW. Theo các nền tảng tổng hợp dữ liệu học thuật được cập nhật năm 2024-2026 (bao gồm các kho lưu trữ Github, Kaggle, và tài liệu từ trường Đại học SINTEF), các mô hình học sâu, hệ thống LNS và hệ thống trí tuệ đa tác tử (multi-agent) như MOGWO-DTACO hoặc DAGA (Dynamic Adaptive Genetic Algorithm) thường đối chiếu hàm mục tiêu tối ưu với "Kỷ lục nghiệm tốt nhất" (BKS - Best Known Solutions). Các BKS này không phải là kết quả của một thuật toán đơn lẻ, mà là một bảng thành tích tổng hợp được đóng góp và tinh chỉnh từng phần nhỏ lẻ từ hàng trăm phương pháp state-of-the-art trong suốt hai thập kỷ qua. Việc một mô hình DRL hay thuật toán metaheuristic như ALNS hoặc HGS thiết lập được một BKS mới trên bất kỳ một instance khó nào của tập dữ liệu Solomon (như nhóm R2 hay RC2) đều được ghi nhận là một đột phá nghiên cứu.Để đánh giá thuật toán ở quy mô tiệm cận thực tế công nghiệp hơn (large-scale problems), bộ dữ liệu mở rộng Gehring & Homberger cung cấp các kịch bản trải dài từ 200, 400, 600, 800 đến tận 1000 khách hàng. Việc giữ vững chất lượng định tuyến và giảm thiểu số lượng xe vận hành trên các hệ thống 1000 nút này đòi hỏi khả năng xử lý song song, thiết lập phân chia cụm siêu việt của mạng thần kinh đồ thị cũng như hệ thống quản lý di truyền khắt khe, loại bỏ hoàn toàn các thuật toán thiết kế kém tối ưu.10.2. Hệ thống cuộc thi thuật toán toàn cầu (DIMACS, GECCO)Thước đo khách quan và uy tín nhất cho sức mạnh giải quyết bài toán VRPTW hiện nay nằm ở kết quả từ các cuộc thi thuật toán toàn cầu. Cuộc thi 12th DIMACS Implementation Challenge dành riêng cho VRP là một trong những nền tảng hội tụ trí tuệ quan trọng nhất của giới vận trù học máy tính, được tổ chức nhằm tôn vinh những thành tựu của David S. Johnson. Dành riêng cho phân nhánh VRPTW, ban tổ chức DIMACS luôn áp đặt các cơ chế đánh giá cực kỳ nghiêm ngặt đối với cửa sổ thời gian cứng (không chấp nhận bất kỳ vi phạm nào và không trao phần thưởng bổ sung cho việc sử dụng ít xe hơn yêu cầu). Cuộc thi được chia làm nhiều giai đoạn đánh giá độc lập (Phase 1, Phase 2) trên các tập dữ liệu bí mật có độ phức tạp khổng lồ để lọc ra những nhóm nghiên cứu sở hữu bộ giải thuật (solvers) tối ưu thực sự thay vì các thuật toán học vẹt trên tập mẫu.Tiếp nối sự dịch chuyển xu hướng ứng dụng điện toán nhận thức sang AI, cuộc thi GECCO 2026 (Tổ chức trên nền tảng ML4VRP) mang chuyên đề "Machine Learning Assisted Evolutionary Computation" đã thiết lập nên một hệ hình thi đấu hoàn toàn mới. Khác với việc nộp các thuật toán được code sẵn, cuộc thi này đòi hỏi những thí sinh phải tận dụng sức mạnh của Deep Learning, Graph Neural Networks (GNNs), Deep Reinforcement Learning (DRL), Neural Combinatorial Optimization, và các Mô hình ngôn ngữ lớn (LLMs) để tiến hóa, hỗ trợ và tự động tạo ra các toán tử định tuyến. Các thí sinh buộc phải khám phá những chân trời kiến trúc mới nhất trên các tập Solomon và Homberger với hàm mục tiêu kép $\text{Objective Function} = 1000 \times NV + TD$ (Trong đó $NV$ là Số lượng phương tiện và $TD$ là Tổng khoảng cách). Điều này minh chứng một điều hiển nhiên: việc định tuyến xe tối ưu không còn là lãnh địa độc quyền của toán học tính toán thuần túy hay thuật toán kinh điển, mà đã vĩnh viễn trở thành điểm hội tụ đỉnh cao của khoa học trí tuệ nhân tạo và học máy tiên tiến nhất.

Lịch sử phát triển của Bài toán Định tuyến Xe có Cửa sổ Thời gian (VRPTW) trong suốt bảy thập kỷ qua là một bức tranh lịch sử sống động, rực rỡ và đầy cảm hứng về sự kiên trì vô hạn trong khoa học Tối ưu hóa Tổ hợp. Khởi nguồn từ việc quy hoạch mô hình toán học sơ khai vài trạm cung cấp xăng dầu với giấy bút vào năm 1959 của Dantzig và Ramser, đến việc thiết lập những nền tảng đánh giá đầu tiên của Solomon những năm 1980, VRPTW đã chứng kiến một quá trình chuyển dịch triết lý giải thuật đầy phức tạp.

Từ sự bảo đảm toán học tối ưu 100% của cấu trúc Branch-and-Cut-and-Price trong việc rẽ nhánh và sinh cột tuyến tính để giải bài toán quy mô nhỏ, sự ra đời của các hệ thống Metaheuristics thích nghi (ALNS, và sau này là tượng đài bất tử Hybrid Genetic Search - HGS) đã bẻ cong giới hạn thời gian tính toán. Các thuật toán này chứng minh rằng việc mạnh dạn đi sâu vào và khám phá không gian bất khả thi (infeasible space) có thể tìm thấy những điểm chạm hoàn hảo để thiết lập Kỷ lục Nghiệm tốt nhất (BKS).

Bước vào thập niên thứ ba của thế kỷ 21, ở ngưỡng cửa năm 2025-2026, VRPTW đã chính thức tiến hóa thành một bài toán được mô hình hóa và giải quyết bằng Mạng thần kinh Đồ thị Chú ý và Học tăng cường sâu (GATv2-DRL, DDQN, Transformer). Sự hỗ trợ không tưởng từ việc tạo code thuật toán tự động thông qua khả năng lý luận (Reason-Act-Reflect) của Mô hình ngôn ngữ lớn LLM (như PyVRP+ hay VRPAgent) và những đột phá phần cứng định tuyến lượng tử (Constraint-aware QAOA, Q-GAT) cùng khả năng tăng tốc tuyến tính 5000x của NVIDIA cuOpt trên GPU đã mở ra một chân trời mới của khái niệm "siêu tự động hóa toán học".

Ở thời điểm hiện tại, vị thế tối ưu tuyệt đối (State-of-the-Art) không còn thuộc về một thuật toán đơn lập nào, mà thuộc về trạng thái lai tạo công nghệ (hybridization). Các hệ thống điều phối logistics công nghiệp mạnh nhất thế giới đang sử dụng năng lực tạo không gian lân cận và khả năng đánh giá tầm nhìn toàn cục của các cấu trúc Transformer hoặc GPU xử lý hàng loạt để bổ trợ trực tiếp cho các bộ gen lai ghép cốt lõi của thuật toán HGS hoặc DRL. Triển vọng tương lai của bài toán VRPTW không chỉ dừng lại ở việc tìm ra đường đi ngắn nhất hay tiết kiệm chi phí một cách tĩnh lặng, mà đang hướng trực tiếp đến khả năng thích ứng động học theo mili-giây tại các môi trường chuỗi cung ứng bất định, mô phỏng sinh thái giao thông phức tạp, điều tiết năng lượng xe điện, và định hình mạng lưới phân phối giao nhận thông minh của tương lai.











\newpage

%========================================================== 



        
\section{Mục tiêu - Phương pháp}
        
Gần đây, phương pháp Adaptive Large Neighborhood Search (ALNS) trở thành tiêu chuẩn nhờ khả năng linh hoạt trong việc phá hủy và tái thiết đồ thị. Dù vậy, ALNS thuần túy bị hạn chế bởi việc chọn toán tử dựa trên xác suất mù quáng. Việc ứng dụng Học tăng cường sâu (Deep Reinforcement Learning) để điều khiển metaheuristics \citep{Author2023} đang mở ra hướng đi mới đầy hứa hẹn, cho phép hệ thống học được các "siêu chiến lược" (meta-policy) từ dữ liệu mô phỏng.

\subsection{Mục tiêu nghiên cứu}
Mục tiêu chính là xây dựng bộ giải \textit{PlateauHybridSolver} sử dụng mạng DDQN nhằm khắc phục tình trạng chững lại của ALNS, đồng thời ngăn chặn hiện tượng thuật toán đánh đổi việc tăng số lượng xe để giảm quãng đường di chuyển.

\subsection{Phương pháp nghiên cứu (Methodology)}

Kiến trúc DDQN-ALNS được mô hình hóa dưới dạng một Quá trình Quyết định Markov (MDP) với ba thành phần cốt lõi: Không gian trạng thái, Không gian hành động, và Hàm phần thưởng. Đặc điểm nổi bật của hệ thống là cơ chế nhận diện điểm chững: DDQN chỉ can thiệp khi ALNS không cải thiện kết quả sau 60 vòng lặp liên tiếp.

\subsubsection{Không gian Trạng thái (State Space)}
Để tác nhân DDQN có thể quan sát cấu trúc và tình trạng tìm kiếm hiện tại mà không bị phụ thuộc vào quy mô đồ thị, một vector trạng thái gồm 10 chiều được trích xuất và chuẩn hóa về khoảng [0, 1]. Phương trình trạng thái $S_t$ được định nghĩa như sau:

\[
S_t = \left[ \text{patience}, \text{gap}, \text{temp\_ratio}, \text{imp\_rate}, \text{nv\_ratio}, \text{route\_balance}, \text{load\_balance}, \text{tw\_tightness}, \text{avg\_slack}, \text{progress} \right]
\]

Các biến số phản ánh độ bế tắc hiện tại (patience), sự chênh lệch so với nghiệm tốt nhất (gap), trạng thái cân bằng tải (load\_balance), và mức độ nghiêm ngặt của các ràng buộc thời gian (tw\_tightness).

\subsubsection{Không gian Hành động (Action Space)}
Thay vì lựa chọn trực tiếp các toán tử cơ sở, DDQN được thiết kế để điều hướng không gian tìm kiếm bằng cách đưa ra 1 trong 4 chiến lược tổng quát (Modes):
\begin{itemize}
    \item \textbf{Mode 0 (Default):} Tiếp tục tỷ lệ xác suất mặc định của ALNS.
    \item \textbf{Mode 1 (Intensify):} Khai thác sâu vùng không gian hiện tại, ưu tiên các toán tử loại bỏ khắt khe như \textit{Worst} hoặc \textit{Shaw removal}.
    \item \textbf{Mode 2 (Diversify):} Phá vỡ cấu trúc nghiệm hiện tại bằng cách tăng cường hệ số phá hủy (destroy\_scale = 1.35) nhằm nhảy khỏi cực trị địa phương.
    \item \textbf{Mode 3 (TW Rescue):} Tập trung xử lý các điểm giao hàng đang vi phạm nghiêm trọng cửa sổ thời gian.
\end{itemize}

\subsubsection{Hàm Phần thưởng nhận thức sức chứa (Capacity-aware Reward)}
Trong tối ưu hóa VRPTW, số lượng xe (NV) đóng vai trò ưu tiên cao hơn tổng quãng đường (TD). Các phương pháp cũ thường gặp "hội chứng trôi dạt", tự ý phân bổ thêm xe để làm giảm TD. Nghiên cứu này đề xuất một hàm phần thưởng $R_t$ áp đặt mức phạt nghiêm khắc (NV Penalty = -15.0) để triệt tiêu hành vi này. Giả sử $\Delta NV = NV_{after} - NV_{before}$ và $\Delta TD = TD_{before} - TD_{after}$ (suy ra $\Delta TD > 0$ là quãng đường đã được rút ngắn), hàm phần thưởng được biểu diễn như sau:

\[
R_t = -0.75 - 0.5 \cdot \mathbb{I}_{(\text{accepted} = 0)} + 
\begin{cases} 
25|\Delta NV| + 100 \max\left(0, \frac{\Delta TD}{TD_{before}}\right) & \text{nếu } \Delta NV < 0 \\ 
-15|\Delta NV| & \text{nếu } \Delta NV > 0 \\ 
250 \left(\frac{\Delta TD}{TD_{before}}\right) & \text{nếu } \Delta NV = 0 \text{ và } \Delta TD > 0 \\ 
0 & \text{ngược lại} 
\end{cases}
\]

(Trong đó $\mathbb{I}_{(\text{accepted} = 0)}$ là hàm chỉ thị, phạt thêm nếu số bước di chuyển bị từ chối).

\section{Kết quả - Thảo luận}
Quá trình huấn luyện trên bộ dữ liệu Solomon RC1 và RC2 chứng minh DDQN-ALNS hoạt động như một \textit{Plateau Controller} hiệu quả. Tại các điểm chững (iteration $> 300$), mạng DDQN liên tục kích hoạt các chiến lược \textit{Diversify} và \textit{Intensify}, giúp chi phí (cost) phá vỡ ngưỡng cực trị địa phương và vượt qua điểm chuẩn (BKS). Đặc biệt, mô hình thể hiện khả năng Zero-shot Transfer Learning xuất sắc khi duy trì mức độ sai số (Gap\%) chỉ 0.70\% trên cấu trúc đồ thị mới lạ hoàn toàn.

\section{Kết luận - Đề nghị}
Công trình đã chứng minh tính hiệu quả của việc kết hợp Học tăng cường sâu như một bộ não điều phối cho Metaheuristics. Hệ thống không chỉ khắc phục được hạn chế kẹt nghiệm mà còn giải quyết triệt để lỗi lạm dụng phương tiện giao hàng. Hướng nghiên cứu tiếp theo sẽ tích hợp Mạng Nơ-ron Đồ thị (GNN) để tự động trích xuất đặc trưng hình học, đồng thời mở rộng thử nghiệm trên các tập dữ liệu có quy mô lớn hơn như Gehring \& Homberger.

% ================= TÀI LIỆU THAM KHẢO =================
\newpage
\begin{thebibliography}{}

\bibitem[Dantzig và Ramser (1959)]{Dantzig1959}
Dantzig, G. B. và Ramser, J. H. (1959). \textit{The Truck Dispatching Problem}. Management Science, 6(1), 80-91.

\bibitem[Solomon (1987)]{Solomon1987}
Solomon, M. M. (1987). \textit{Algorithms for the Vehicle Routing and Scheduling Problems with Time Window Constraints}. Operations Research, 35(2), 254-265.

\end{thebibliography}

\end{document}

Giả sử mạng lưới giao thông được biểu diễn bằng một đồ thị có hướng $G = (V, A)$, trong đó $V = \{0, 1,..., n\}$ là tập hợp các đỉnh (nút). Nút $0$ đại diện cho trạm lưu trữ trung tâm (depot), nơi xuất phát và kết thúc của mọi phương tiện, trong khi tập $N = V \setminus \{0\}$ đại diện cho tập hợp $n$ khách hàng phân tán về mặt địa lý.\\ Tập hợp $A = \{(i, j) \mid i, j \in V, i \neq j\}$ đại diện cho các cung đường kết nối giữa các nút trong mạng lưới.


Giả sử mạng lưới giao thông được biểu diễn bằng một đồ thị có hướng $G = (V, A)$, trong đó $V = \{0, 1,..., n\}$ là tập hợp các đỉnh (nút). Nút $0$ đại diện cho trạm lưu trữ trung tâm (depot), nơi xuất phát và kết thúc của mọi phương tiện, trong khi tập $N = V \setminus \{0\}$ đại diện cho tập hợp $n$ khách hàng phân tán về mặt địa lý.\\ Tập hợp $A = \{(i, j) \mid i, j \in V, i \neq j\}$ đại diện cho các cung đường kết nối giữa các nút trong mạng lưới.

Mỗi khách hàng $i \in N$ có các tham số liên kết nội tại như sau:Nhu cầu tài nguyên (hàng hóa): $q_i \geq 0$.Cửa sổ thời gian: $[e_i, l_i]$, trong đó $e_i$ là thời gian bắt đầu phục vụ sớm nhất có thể và $l_i$ là thời gian trễ nhất cho phép.Thời gian phục vụ (service time): $s_i \geq 0$, là khoảng thời gian cần thiết để thực hiện thao tác dỡ hoặc tải hàng tại địa điểm của khách hàng $i$.Mỗi cung đường $(i, j) \in A$ gắn liền với:Chi phí di chuyển (routing cost / distance): $c_{ij}$.Thời gian di chuyển (travel time): $t_{ij}$, trong nhiều trường hợp nghiên cứu chuẩn hóa, thời gian di chuyển được giả định tỷ lệ thuận hoặc bằng với khoảng cách di chuyển.

Hệ thống điều phối sở hữu một đội xe bao gồm $K$ phương tiện đồng nhất, mỗi phương tiện có sức chứa vật lý tối đa là $Q$. Các biến quyết định hình thành nên không gian tìm kiếm của bài toán bao gồm:$x_{ij}^k \in \{0, 1\}$: Biến nhị phân nhận giá trị $1$ nếu phương tiện $k$ di chuyển trực tiếp từ nút $i$ đến nút $j$, và nhận giá trị $0$ nếu ngược lại.$w_i^k \geq 0$: Biến liên tục biểu thị thời điểm thực tế bắt đầu phục vụ khách hàng $i$ bởi phương tiện $k$.Hàm mục tiêu tiêu chuẩn của VRPTW thường mang tính phân cấp (hierarchical objective structure), phản ánh ưu tiên kinh tế của các doanh nghiệp logistics: Ưu tiên hàng đầu (Primary Objective) là giảm thiểu tổng số phương tiện sử dụng (tương đương với chi phí cố định cho nhân sự và khấu hao xe), sau đó ưu tiên thứ hai (Secondary Objective) là giảm thiểu tổng chi phí hoặc tổng khoảng cách di chuyển (tương đương với chi phí nhiên liệu).
'Hàm mục tiêu có thể được biểu diễn như sau:Hàm mục tiêu 1 (Số lượng xe): $\min \sum_{k=1}^K \sum_{j \in N} x_{0j}^k$Hàm mục tiêu 2 (Quãng đường): $\min \sum_{k=1}^K \sum_{(i, j) \in A} c_{ij} x_{ij}^k$Các hệ ràng buộc cơ bản làm nên tính chất khó của VRPTW bao gồm:Ràng buộc phân công phục vụ (Assignment constraint): Đảm bảo mỗi khách hàng được viếng thăm chính xác một lần bởi một phương tiện duy nhất trong toàn bộ hệ thống tuyến đường:$\sum_{k=1}^K \sum_{j \in V, j \neq i} x_{ij}^k = 1, \quad \forall i \in N$Ràng buộc bảo toàn luồng (Flow conservation constraint): Tính liên tục của định tuyến; một phương tiện khi đi vào một nút khách hàng thì cũng bắt buộc phải đi ra từ nút đó để tiếp tục hành trình:$\sum_{i \in V} x_{ih}^k - \sum_{j \in V} x_{hj}^k = 0, \quad \forall h \in N, \forall k$Ràng buộc luồng tại trạm gốc (Depot constraints): Mọi tuyến đường đang hoạt động phải có điểm khởi đầu và điểm kết thúc tại trạm gốc:$\sum_{j \in N} x_{0j}^k \leq 1, \quad \forall k$$\sum_{i \in N} x_{i0}^k \leq 1, \quad \forall k$Ràng buộc sức chứa tải trọng (Capacity constraint): Tổng nhu cầu của tất cả các khách hàng được phục vụ trên một tuyến đường không bao giờ được vượt quá sức chứa tối đa $Q$ của phương tiện đó :
$\sum_{i \in N} q_i \sum_{j \in V} x_{ij}^k \leq Q, \quad \forall k$Ràng buộc liên tục lịch trình thời gian (Time schedule constraint): Đảm bảo tính nhất quán của thời gian dọc theo lộ trình. Nếu xe đi từ $i$ đến $j$, thời gian đến $j$ phải bằng hoặc lớn hơn thời gian bắt đầu tại $i$ cộng với thời gian phục vụ tại $i$ và thời gian di chuyển từ $i$ đến $j$. Ràng buộc này thường được thiết lập bằng phương pháp Big-M:$w_i^k + s_i + t_{ij} - M(1 - x_{ij}^k) \leq w_j^k, \quad \forall (i, j) \in A, \forall k$Trong đó $M$ là một hằng số dương đủ lớn để vô hiệu hóa ràng buộc khi $x_{ij}^k = 0$.Ràng buộc cửa sổ thời gian (Time window bounds): Việc bắt đầu phục vụ phải tuân thủ nghiêm ngặt ranh giới đã thỏa thuận với khách hàng:
$e_i \leq w_i^k \leq l_i, \quad \forall i \in V, \forall k$.2.1. Phân loại Cửa sổ thời gian: Cứng và MềmMột khía cạnh quan trọng trong mô hình hóa VRPTW là sự phân định giữa Cửa sổ Thời gian Cứng (Hard Time Windows - VRPHTW) và Cửa sổ Thời gian Mềm (Soft Time Windows - VRPSTW).
Trong các mô hình cứng (được mô tả bằng công thức trên), các giới hạn thời gian tuyệt đối không được phép vi phạm. Nếu một phương tiện đến sớm, nó phải đợi; nếu nó dự kiến đến muộn, giải pháp đó lập tức bị phân loại là bất khả thi (infeasible) và bị loại bỏ khỏi không gian tìm kiếm. Các cuộc thi lớn như DIMACS Implementation Challenge thường áp dụng quy ước cửa sổ thời gian cứng.
Ngược lại, cửa sổ thời gian mềm cho phép phương tiện đến sớm hơn hoặc muộn hơn khoảng thời gian ưu tiên của khách hàng, nhưng mọi sự vi phạm đều bị tính một khoản chi phí phạt (penalty cost) vào hàm mục tiêu. Điều này tạo ra một sự linh hoạt nhất định, giúp các thuật toán dễ dàng tìm thấy các giải pháp hơn, nhưng lại đòi hỏi cơ chế cân bằng chi phí cực kỳ tinh tế. Báo cáo này chủ yếu tập trung vào tiến trình của mô hình cửa sổ cứng, do nó đại diện cho thách thức toán học khốc liệt nhất.
